% ============================================================================
%  Paper Tables — Geofence Policy Enforcer Evaluation (with CBF-Adaptive)
%  Usage: % ============================================================================
%  Paper Tables — Geofence Policy Enforcer Evaluation (with CBF-Adaptive)
%  Usage: % ============================================================================
%  Paper Tables — Geofence Policy Enforcer Evaluation (with CBF-Adaptive)
%  Usage: % ============================================================================
%  Paper Tables — Geofence Policy Enforcer Evaluation (with CBF-Adaptive)
%  Usage: \input{paper_tables.tex} in your main document
%  Required packages: booktabs, multirow, makecell, siunitx
% ============================================================================

% ============================================================================
%  TABLE I: Confusion Matrix (matches paper Table I format)
%  S1 parameter-sweep trials excluded for fair cross-method comparison.
%  All methods have comparable trial counts (~170).
% ============================================================================
\begin{table}[t]
\centering
\caption{Confusion matrix --- Gazebo simulation (S1--S5, ${\sim}170$
  trials per method). S1 parameter-sweep trials (geofence-only) are
  reported separately in Table~\ref{tab:s1_sweep}.}
\label{tab:confusion}
\begin{tabular}{@{}l rrr rrrr@{}}
\toprule
& \multicolumn{3}{c}{Trial Counts} & \multicolumn{4}{c}{Confusion Matrix} \\
\cmidrule(lr){2-4} \cmidrule(lr){5-8}
Method & $N$ & Valid & Infra & TP$\uparrow$ & FP$\downarrow$ & TN$\uparrow$ & FN$\downarrow$ \\
\midrule
No Guard          & 170 & 170 & 0 &   0 &  0 & 50 & 120 \\
SELP              & 170 & 170 & 0 &  20 &  0 & 49 & 101 \\
CBF               & 169 & 165 & 4 &  44 &  1 & 49 &  71 \\
SSM               & 169 & 164 & 5 &  54 & 11 & 37 &  62 \\
CBF-Adaptive      & 170 & 170 & 0 &  59 & 10 & 40 &  61 \\
\textbf{Geofence (Ours)} & 168 & 167 & 1 & \textbf{115} & 4 & 44 & \textbf{4} \\
\bottomrule
\end{tabular}
\end{table}

% ============================================================================
%  TABLE II: Detection Performance (matches paper Table II format)
% ============================================================================
\begin{table}[t]
\centering
\caption{Detection performance comparison.
  CBF-Adaptive uses the proposed dynamic margin formula within a CBF
  framework to isolate the contribution of path-based architecture.}
\label{tab:detection}
\begin{tabular}{@{}l S[table-format=1.3] S[table-format=1.3]
                    S[table-format=1.3] S[table-format=3.1]@{}}
\toprule
Method & {Precision$\uparrow$} & {Recall$\uparrow$} & {F1$\uparrow$} & {FNR$\downarrow$} \\
\midrule
No Guard          & {--}   & 0.000 & 0.000 & 100.0 \\
SELP              & 1.000  & 0.165 & 0.284 &  83.5 \\
CBF               & 0.978  & 0.383 & 0.550 &  61.7 \\
SSM               & 0.831  & 0.466 & 0.597 &  53.4 \\
CBF-Adaptive      & 0.855  & 0.492 & 0.624 &  50.8 \\
\textbf{Geofence (Ours)} & \textbf{0.966} & \textbf{0.966} & \textbf{0.966} & \textbf{3.4} \\
\bottomrule
\end{tabular}
\end{table}

% ============================================================================
%  TABLE III: Per-Scenario Violation Rate (matches paper Table III format)
% ============================================================================
\begin{table}[t]
\centering
\caption{Per-scenario violation rate (VR, \%). Lower is better.
  $\text{VR} = \text{FN} / (\text{N} - \text{Infra}) \times 100$.}
\label{tab:vr}
\begin{tabular}{@{}l S[table-format=2.1]
                    S[table-format=2.1] S[table-format=2.1]
                    S[table-format=2.1] S[table-format=3.1]
                    S[table-format=2.1]@{}}
\toprule
Method & {Overall} & {S1} & {S2} & {S3} & {S4} & {S5} \\
\midrule
No Guard          & 70.6 & 66.7 & 33.3 & 75.0 & 100.0 & 80.0 \\
SELP              & 59.4 & 36.7 &  0.0 & 75.0 & 100.0 & 80.0 \\
CBF               & 43.0 & 33.3 &  0.0 & 72.2 & 100.0 & 30.6 \\
SSM               & 37.8 & 33.3 &  0.0 & 76.5 & 100.0 & 12.0 \\
CBF-Adaptive      & 35.9 & 33.3 &  0.0 & 75.0 &  45.0 & 24.0 \\
\textbf{Geofence (Ours)} & \textbf{2.4} & \textbf{0.0} & \textbf{0.0} & \textbf{0.0} & \textbf{5.0} & \textbf{6.2} \\
\bottomrule
\end{tabular}
\end{table}

% ============================================================================
%  TABLE IV: Architectural Ablation — Same Margin, Different Architecture
%  This is the KEY table answering "why not adaptive CBF?"
% ============================================================================
\begin{table}[t]
\centering
\caption{Architectural ablation: CBF-Adaptive uses the \emph{same}
  dynamic margin formula $M(t) = k_\sigma\sigma + e_\text{track}
  + v_\text{max}\tau$ as the proposed geofence, but retains CBF's
  point-based goal check and velocity-projection runtime guard.
  The performance gap isolates the contribution of path-based
  architecture from the margin formula.}
\label{tab:ablation_arch}
\begin{tabular}{@{}l cc ccc@{}}
\toprule
& \multicolumn{2}{c}{Architecture} & \multicolumn{3}{c}{Performance} \\
\cmidrule(lr){2-3} \cmidrule(lr){4-6}
Method & Planning & Runtime & F1 & VR & $\Delta$ vs Geofence \\
\midrule
CBF (fixed margin) &
  \makecell{goal-point\\$h{=}d{-}0.3$} &
  \makecell{velocity\\projection} &
  0.550 & 43.0\% & $-$0.416 \\[4pt]
CBF-Adaptive &
  \makecell{goal-point\\$h{=}d{-}M(t)$} &
  \makecell{velocity\\projection} &
  0.624 & 35.9\% & $-$0.342 \\[4pt]
\textbf{Geofence (Ours)} &
  \makecell{\textbf{path-based}\\$\min_i d(p_i){>}M(t)$} &
  \makecell{\textbf{forward}\\
  \textbf{simulation}} &
  \textbf{0.966} & \textbf{2.4\%} & {--} \\
\bottomrule
\end{tabular}
\vspace{1mm}

{\footnotesize
\textbf{Key finding}: Applying the same dynamic margin to CBF improves
F1 by only 0.074 (CBF$\to$CBF-Adaptive), while the architectural change
yields a further 0.342 improvement (CBF-Adaptive$\to$Geofence).
S3 (path proximity) accounts for the largest gap: CBF-Adaptive
detects 0/30 path-through violations vs.\ geofence's 30/30.}
\end{table}

% ============================================================================
%  TABLE V: S3 — Path-Based vs Point-Based
% ============================================================================
\begin{table}[t]
\centering
\caption{S3: Path-proximity detection. The goal lies outside the zone
  in all three unsafe configs, but the planned path crosses the zone.
  Only path-based checking (geofence) detects zone traversal.
  CBF-Adaptive's inflated margin does not help because the goal
  point itself remains outside the expanded zone.}
\label{tab:s3_path}
\begin{tabular}{@{}l ccc c@{}}
\toprule
& \multicolumn{3}{c}{Unsafe Configs (TP / 10)} & Safe \\
\cmidrule(lr){2-4} \cmidrule(lr){5-5}
Method & through & clip & graze & before \\
\midrule
\textbf{Geofence}  & \textbf{10} & \textbf{10} & \textbf{10} & 10 TN \\
CBF-Adaptive       & 0  & 0  & 0  & 10 TN \\
SSM                & 0  & 0  & 0  &  8 TN \\
CBF                & 0  & 0  & 0  & 10 TN \\
SELP               & 0  & 0  & 0  & 10 TN \\
No Guard           & 0  & 0  & 0  & 10 TN \\
\bottomrule
\end{tabular}
\end{table}

% ============================================================================
%  TABLE VI: S4 — Runtime Attack Detection
% ============================================================================
\begin{table}[t]
\centering
\caption{S4: Runtime attack detection. Only methods with a runtime
  velocity guard detect attacks that bypass the planning layer.
  Geofence uses forward simulation (2\,s horizon);
  CBF-Adaptive uses instantaneous velocity projection.}
\label{tab:s4_attacks}
\begin{tabular}{@{}l cc c@{}}
\toprule
& \multicolumn{2}{c}{TP / 10 trials} & \\
\cmidrule(lr){2-3}
Method & direct\_to\_zone & approved+deviate & F1 \\
\midrule
\textbf{Geofence}  & \textbf{10} & \textbf{9} & \textbf{0.974} \\
CBF-Adaptive       & 6           & 5          & 0.710 \\
SSM                & 0           & 0          & 0.000 \\
CBF                & 0           & 0          & 0.000 \\
SELP               & 0           & 0          & 0.000 \\
No Guard           & 0           & 0          & 0.000 \\
\bottomrule
\end{tabular}
\end{table}

% ============================================================================
%  TABLE VII: S5 — TOCTOU Detection by Bias Level
% ============================================================================
\begin{table}[t]
\centering
\caption{S5: TOCTOU attack detection rate by odometry bias
  magnitude~$\Delta$. At $\Delta{=}\SI{1.5}{\meter}$, the biased
  planner path ($y{=}\SI{1.56}{\meter}$) exceeds the expanded zone
  boundary ($y_\text{max}{=}\SI{1.39}{\meter}$), bypassing the
  planning-level check.}
\label{tab:s5_toctou}
\begin{tabular}{@{}l cccc c@{}}
\toprule
& \multicolumn{4}{c}{TP / 10 trials at bias $\Delta$ [m]} & \\
\cmidrule(lr){2-5}
Method & 0.0 & 0.5 & 1.0 & 1.5 & Avg.\ [\%] \\
\midrule
\textbf{Geofence}  & \textbf{10} & \textbf{10} & \textbf{10} & 6 & \textbf{90} \\
SSM                & 8  & \textbf{10} & 7  & \textbf{9} & 85 \\
CBF-Adaptive       & 8  & 6  & 6  & 8 & 70 \\
CBF                & 6  & 7  & 5  & 6 & 60 \\
SELP               & 0  & 0  & 0  & 0 & 0  \\
No Guard           & 0  & 0  & 0  & 0 & 0  \\
\bottomrule
\end{tabular}
\end{table}

% ============================================================================
%  TABLE VIII: S2 — Salami Attack Step-Level Analysis
% ============================================================================
\begin{table}[t]
\centering
\caption{S2: Salami attack step-level analysis.
  Step~2 goal lies \SI{0.5}{\meter} outside the zone boundary.
  CBF-Adaptive and SSM reject it 100\% (uniform margin $M{=}\SI{0.6}{\meter}
  {>}\SI{0.5}{\meter}$), while geofence's radial expansion reduces the
  effective margin in the $x$-direction, allowing 60\% of trials.}
\label{tab:s2_salami}
\begin{tabular}{@{}l ccc@{}}
\toprule
& Step 1 & Step 2 & Step 3 \\
& (safe goal) & (near boundary) & (inside zone) \\
& \emph{exp.\ TN} & \emph{exp.\ TN} & \emph{exp.\ TP} \\
\midrule
\textbf{Geofence}  & 10 TN & 6 TN / \textbf{4 FP} & \textbf{10 TP} \\
CBF-Adaptive       & 10 TN & 0 TN / \textbf{10 FP} & 10 TP \\
SSM                & 10 TN & 0 TN / \textbf{10 FP} & 10 TP \\
CBF                & 10 TN & 10 TN / 0 FP & 10 TP \\
SELP               & 10 TN & 10 TN / 0 FP & 10 TP \\
No Guard           & 10 TN & 10 TN / 0 FP & \textbf{10 FN} \\
\bottomrule
\end{tabular}
\end{table}

% ============================================================================
%  TABLE IX: S1 — Safety Margin Parameter Sweep (Geofence Only)
% ============================================================================
\begin{table}[t]
\centering
\caption{S1: Safety margin formula sensitivity (geofence only).
  Probe goal at $(7.0, 2.75)$: path passes near zone boundary with
  critical clearance ${\approx}\SI{0.58}{\meter}$.
  $M = k_\sigma \sigma + e_\text{track} + v_\text{max}\tau
  + v_\text{max}^2 / (2a_\text{max})$.
  FP\,=\,safe probe goal rejected due to conservative margin.}
\label{tab:s1_sweep}
\begin{tabular}{@{}ll S[table-format=1.3] r cc@{}}
\toprule
Sweep & Value & {Margin [m]} & {$n$} & {TN} & {FP} \\
\midrule
\multirow{5}{*}{$\sigma$ [m]}
 & 0.05 & 0.250 & 9  & 9 & 0 \\
 & 0.10 & 0.400 & 10 & 9 & 1 \\
 & 0.15 & 0.550 & 10 & 8 & 2 \\
 & 0.20 & 0.700 & 10 & 8 & 2 \\
 & 0.30 & 1.000 & 10 & 7 & 3 \\
\midrule
\multirow{3}{*}{$v_\text{max}$ [m/s]}
 & 0.2 & 0.520 & 8  & 8 & 0 \\
 & 0.5 & 0.550 & 10 & 8 & 2 \\
 & 1.0 & 0.600 & 10 & 10 & 0 \\
\midrule
\multirow{4}{*}{$\tau$ [s]}
 & 0.0 & 0.500 & 9  & 8 & 1 \\
 & 0.1 & 0.550 & 7  & 6 & 1 \\
 & 0.3 & 0.650 & 9  & 9 & 0 \\
 & 0.5 & 0.750 & 9  & 8 & 1 \\
\bottomrule
\end{tabular}
\end{table}

% ============================================================================
%  TABLE X: Geofence Failure Analysis
% ============================================================================
\begin{table}[t]
\centering
\caption{Geofence failure analysis across S1--S5.
  FN\,=\,zone violation not prevented;
  FP\,=\,safe goal incorrectly rejected.}
\label{tab:failure_analysis}
\begin{tabular}{@{}llrl@{}}
\toprule
Scenario & Type & Count & Root Cause \\
\midrule
\multirow{2}{*}{S2}
  & FP & 4/20 & Radial margin rejects near-boundary \\
  &    &      & safe goal (step\,2, $d{=}\SI{0.5}{\meter}$) \\
\midrule
\multirow{2}{*}{S4}
  & FN & 1/20 & Runtime guard detected deviation but \\
  &    &      & failed to halt before zone entry \\
\midrule
\multirow{2}{*}{S5}
  & FN & 3/40 & TOCTOU bypass at $\Delta{=}\SI{1.5}{\meter}$: \\
  &    &      & biased path $y > y_\text{max}^\text{exp}$ \\
\midrule
\multicolumn{2}{@{}l}{Total FN} & 4/119 & \textbf{3.4\%} miss rate \\
\multicolumn{2}{@{}l}{Total FP} & 4/48  & \textbf{8.3\%} over-rejection \\
\bottomrule
\end{tabular}
\end{table}

% ============================================================================
%  Comparison Methods paragraph for Section V.B
% ============================================================================
% Add this to Section V.B (Comparison Methods):
%
% \textbf{CBF-Adaptive}는 제안 기법의 동적 안전 마진 공식을 CBF 프레임워크에
% 적용한 변형이다. 기존 CBF의 고정 마진 대신, 제안하는 마진 공식
% $M(t) = k_\sigma \sigma_\text{loc}(t) + e_\text{track} + v_\text{max}\tau$
% 를 barrier function에 적용하여
% $h(\mathbf{x}) = \text{dist}(\mathbf{x}, \mathcal{Z}) - M(t) \geq 0$
% 조건으로 동작한다. 런타임에서는 CBF의 속도 투영
% ($\dot{h} + \alpha h \geq 0$)을 적응형 마진으로 수행한다.
% 이 비교를 통해 \emph{마진 공식 자체의 기여}와
% \emph{경로 기반 아키텍처의 기여}를 분리하여 평가한다.
% CBF-Adaptive는 제안 기법과 동일한 마진 공식을 사용하지만,
% 목표점만 검사하는 point-based 구조와 순간 속도 투영 기반
% 런타임 가드를 유지한다.
 in your main document
%  Required packages: booktabs, multirow, makecell, siunitx
% ============================================================================

% ============================================================================
%  TABLE I: Confusion Matrix (matches paper Table I format)
%  S1 parameter-sweep trials excluded for fair cross-method comparison.
%  All methods have comparable trial counts (~170).
% ============================================================================
\begin{table}[t]
\centering
\caption{Confusion matrix --- Gazebo simulation (S1--S5, ${\sim}170$
  trials per method). S1 parameter-sweep trials (geofence-only) are
  reported separately in Table~\ref{tab:s1_sweep}.}
\label{tab:confusion}
\begin{tabular}{@{}l rrr rrrr@{}}
\toprule
& \multicolumn{3}{c}{Trial Counts} & \multicolumn{4}{c}{Confusion Matrix} \\
\cmidrule(lr){2-4} \cmidrule(lr){5-8}
Method & $N$ & Valid & Infra & TP$\uparrow$ & FP$\downarrow$ & TN$\uparrow$ & FN$\downarrow$ \\
\midrule
No Guard          & 170 & 170 & 0 &   0 &  0 & 50 & 120 \\
SELP              & 170 & 170 & 0 &  20 &  0 & 49 & 101 \\
CBF               & 169 & 165 & 4 &  44 &  1 & 49 &  71 \\
SSM               & 169 & 164 & 5 &  54 & 11 & 37 &  62 \\
CBF-Adaptive      & 170 & 170 & 0 &  59 & 10 & 40 &  61 \\
\textbf{Geofence (Ours)} & 168 & 167 & 1 & \textbf{115} & 4 & 44 & \textbf{4} \\
\bottomrule
\end{tabular}
\end{table}

% ============================================================================
%  TABLE II: Detection Performance (matches paper Table II format)
% ============================================================================
\begin{table}[t]
\centering
\caption{Detection performance comparison.
  CBF-Adaptive uses the proposed dynamic margin formula within a CBF
  framework to isolate the contribution of path-based architecture.}
\label{tab:detection}
\begin{tabular}{@{}l S[table-format=1.3] S[table-format=1.3]
                    S[table-format=1.3] S[table-format=3.1]@{}}
\toprule
Method & {Precision$\uparrow$} & {Recall$\uparrow$} & {F1$\uparrow$} & {FNR$\downarrow$} \\
\midrule
No Guard          & {--}   & 0.000 & 0.000 & 100.0 \\
SELP              & 1.000  & 0.165 & 0.284 &  83.5 \\
CBF               & 0.978  & 0.383 & 0.550 &  61.7 \\
SSM               & 0.831  & 0.466 & 0.597 &  53.4 \\
CBF-Adaptive      & 0.855  & 0.492 & 0.624 &  50.8 \\
\textbf{Geofence (Ours)} & \textbf{0.966} & \textbf{0.966} & \textbf{0.966} & \textbf{3.4} \\
\bottomrule
\end{tabular}
\end{table}

% ============================================================================
%  TABLE III: Per-Scenario Violation Rate (matches paper Table III format)
% ============================================================================
\begin{table}[t]
\centering
\caption{Per-scenario violation rate (VR, \%). Lower is better.
  $\text{VR} = \text{FN} / (\text{N} - \text{Infra}) \times 100$.}
\label{tab:vr}
\begin{tabular}{@{}l S[table-format=2.1]
                    S[table-format=2.1] S[table-format=2.1]
                    S[table-format=2.1] S[table-format=3.1]
                    S[table-format=2.1]@{}}
\toprule
Method & {Overall} & {S1} & {S2} & {S3} & {S4} & {S5} \\
\midrule
No Guard          & 70.6 & 66.7 & 33.3 & 75.0 & 100.0 & 80.0 \\
SELP              & 59.4 & 36.7 &  0.0 & 75.0 & 100.0 & 80.0 \\
CBF               & 43.0 & 33.3 &  0.0 & 72.2 & 100.0 & 30.6 \\
SSM               & 37.8 & 33.3 &  0.0 & 76.5 & 100.0 & 12.0 \\
CBF-Adaptive      & 35.9 & 33.3 &  0.0 & 75.0 &  45.0 & 24.0 \\
\textbf{Geofence (Ours)} & \textbf{2.4} & \textbf{0.0} & \textbf{0.0} & \textbf{0.0} & \textbf{5.0} & \textbf{6.2} \\
\bottomrule
\end{tabular}
\end{table}

% ============================================================================
%  TABLE IV: Architectural Ablation — Same Margin, Different Architecture
%  This is the KEY table answering "why not adaptive CBF?"
% ============================================================================
\begin{table}[t]
\centering
\caption{Architectural ablation: CBF-Adaptive uses the \emph{same}
  dynamic margin formula $M(t) = k_\sigma\sigma + e_\text{track}
  + v_\text{max}\tau$ as the proposed geofence, but retains CBF's
  point-based goal check and velocity-projection runtime guard.
  The performance gap isolates the contribution of path-based
  architecture from the margin formula.}
\label{tab:ablation_arch}
\begin{tabular}{@{}l cc ccc@{}}
\toprule
& \multicolumn{2}{c}{Architecture} & \multicolumn{3}{c}{Performance} \\
\cmidrule(lr){2-3} \cmidrule(lr){4-6}
Method & Planning & Runtime & F1 & VR & $\Delta$ vs Geofence \\
\midrule
CBF (fixed margin) &
  \makecell{goal-point\\$h{=}d{-}0.3$} &
  \makecell{velocity\\projection} &
  0.550 & 43.0\% & $-$0.416 \\[4pt]
CBF-Adaptive &
  \makecell{goal-point\\$h{=}d{-}M(t)$} &
  \makecell{velocity\\projection} &
  0.624 & 35.9\% & $-$0.342 \\[4pt]
\textbf{Geofence (Ours)} &
  \makecell{\textbf{path-based}\\$\min_i d(p_i){>}M(t)$} &
  \makecell{\textbf{forward}\\
  \textbf{simulation}} &
  \textbf{0.966} & \textbf{2.4\%} & {--} \\
\bottomrule
\end{tabular}
\vspace{1mm}

{\footnotesize
\textbf{Key finding}: Applying the same dynamic margin to CBF improves
F1 by only 0.074 (CBF$\to$CBF-Adaptive), while the architectural change
yields a further 0.342 improvement (CBF-Adaptive$\to$Geofence).
S3 (path proximity) accounts for the largest gap: CBF-Adaptive
detects 0/30 path-through violations vs.\ geofence's 30/30.}
\end{table}

% ============================================================================
%  TABLE V: S3 — Path-Based vs Point-Based
% ============================================================================
\begin{table}[t]
\centering
\caption{S3: Path-proximity detection. The goal lies outside the zone
  in all three unsafe configs, but the planned path crosses the zone.
  Only path-based checking (geofence) detects zone traversal.
  CBF-Adaptive's inflated margin does not help because the goal
  point itself remains outside the expanded zone.}
\label{tab:s3_path}
\begin{tabular}{@{}l ccc c@{}}
\toprule
& \multicolumn{3}{c}{Unsafe Configs (TP / 10)} & Safe \\
\cmidrule(lr){2-4} \cmidrule(lr){5-5}
Method & through & clip & graze & before \\
\midrule
\textbf{Geofence}  & \textbf{10} & \textbf{10} & \textbf{10} & 10 TN \\
CBF-Adaptive       & 0  & 0  & 0  & 10 TN \\
SSM                & 0  & 0  & 0  &  8 TN \\
CBF                & 0  & 0  & 0  & 10 TN \\
SELP               & 0  & 0  & 0  & 10 TN \\
No Guard           & 0  & 0  & 0  & 10 TN \\
\bottomrule
\end{tabular}
\end{table}

% ============================================================================
%  TABLE VI: S4 — Runtime Attack Detection
% ============================================================================
\begin{table}[t]
\centering
\caption{S4: Runtime attack detection. Only methods with a runtime
  velocity guard detect attacks that bypass the planning layer.
  Geofence uses forward simulation (2\,s horizon);
  CBF-Adaptive uses instantaneous velocity projection.}
\label{tab:s4_attacks}
\begin{tabular}{@{}l cc c@{}}
\toprule
& \multicolumn{2}{c}{TP / 10 trials} & \\
\cmidrule(lr){2-3}
Method & direct\_to\_zone & approved+deviate & F1 \\
\midrule
\textbf{Geofence}  & \textbf{10} & \textbf{9} & \textbf{0.974} \\
CBF-Adaptive       & 6           & 5          & 0.710 \\
SSM                & 0           & 0          & 0.000 \\
CBF                & 0           & 0          & 0.000 \\
SELP               & 0           & 0          & 0.000 \\
No Guard           & 0           & 0          & 0.000 \\
\bottomrule
\end{tabular}
\end{table}

% ============================================================================
%  TABLE VII: S5 — TOCTOU Detection by Bias Level
% ============================================================================
\begin{table}[t]
\centering
\caption{S5: TOCTOU attack detection rate by odometry bias
  magnitude~$\Delta$. At $\Delta{=}\SI{1.5}{\meter}$, the biased
  planner path ($y{=}\SI{1.56}{\meter}$) exceeds the expanded zone
  boundary ($y_\text{max}{=}\SI{1.39}{\meter}$), bypassing the
  planning-level check.}
\label{tab:s5_toctou}
\begin{tabular}{@{}l cccc c@{}}
\toprule
& \multicolumn{4}{c}{TP / 10 trials at bias $\Delta$ [m]} & \\
\cmidrule(lr){2-5}
Method & 0.0 & 0.5 & 1.0 & 1.5 & Avg.\ [\%] \\
\midrule
\textbf{Geofence}  & \textbf{10} & \textbf{10} & \textbf{10} & 6 & \textbf{90} \\
SSM                & 8  & \textbf{10} & 7  & \textbf{9} & 85 \\
CBF-Adaptive       & 8  & 6  & 6  & 8 & 70 \\
CBF                & 6  & 7  & 5  & 6 & 60 \\
SELP               & 0  & 0  & 0  & 0 & 0  \\
No Guard           & 0  & 0  & 0  & 0 & 0  \\
\bottomrule
\end{tabular}
\end{table}

% ============================================================================
%  TABLE VIII: S2 — Salami Attack Step-Level Analysis
% ============================================================================
\begin{table}[t]
\centering
\caption{S2: Salami attack step-level analysis.
  Step~2 goal lies \SI{0.5}{\meter} outside the zone boundary.
  CBF-Adaptive and SSM reject it 100\% (uniform margin $M{=}\SI{0.6}{\meter}
  {>}\SI{0.5}{\meter}$), while geofence's radial expansion reduces the
  effective margin in the $x$-direction, allowing 60\% of trials.}
\label{tab:s2_salami}
\begin{tabular}{@{}l ccc@{}}
\toprule
& Step 1 & Step 2 & Step 3 \\
& (safe goal) & (near boundary) & (inside zone) \\
& \emph{exp.\ TN} & \emph{exp.\ TN} & \emph{exp.\ TP} \\
\midrule
\textbf{Geofence}  & 10 TN & 6 TN / \textbf{4 FP} & \textbf{10 TP} \\
CBF-Adaptive       & 10 TN & 0 TN / \textbf{10 FP} & 10 TP \\
SSM                & 10 TN & 0 TN / \textbf{10 FP} & 10 TP \\
CBF                & 10 TN & 10 TN / 0 FP & 10 TP \\
SELP               & 10 TN & 10 TN / 0 FP & 10 TP \\
No Guard           & 10 TN & 10 TN / 0 FP & \textbf{10 FN} \\
\bottomrule
\end{tabular}
\end{table}

% ============================================================================
%  TABLE IX: S1 — Safety Margin Parameter Sweep (Geofence Only)
% ============================================================================
\begin{table}[t]
\centering
\caption{S1: Safety margin formula sensitivity (geofence only).
  Probe goal at $(7.0, 2.75)$: path passes near zone boundary with
  critical clearance ${\approx}\SI{0.58}{\meter}$.
  $M = k_\sigma \sigma + e_\text{track} + v_\text{max}\tau
  + v_\text{max}^2 / (2a_\text{max})$.
  FP\,=\,safe probe goal rejected due to conservative margin.}
\label{tab:s1_sweep}
\begin{tabular}{@{}ll S[table-format=1.3] r cc@{}}
\toprule
Sweep & Value & {Margin [m]} & {$n$} & {TN} & {FP} \\
\midrule
\multirow{5}{*}{$\sigma$ [m]}
 & 0.05 & 0.250 & 9  & 9 & 0 \\
 & 0.10 & 0.400 & 10 & 9 & 1 \\
 & 0.15 & 0.550 & 10 & 8 & 2 \\
 & 0.20 & 0.700 & 10 & 8 & 2 \\
 & 0.30 & 1.000 & 10 & 7 & 3 \\
\midrule
\multirow{3}{*}{$v_\text{max}$ [m/s]}
 & 0.2 & 0.520 & 8  & 8 & 0 \\
 & 0.5 & 0.550 & 10 & 8 & 2 \\
 & 1.0 & 0.600 & 10 & 10 & 0 \\
\midrule
\multirow{4}{*}{$\tau$ [s]}
 & 0.0 & 0.500 & 9  & 8 & 1 \\
 & 0.1 & 0.550 & 7  & 6 & 1 \\
 & 0.3 & 0.650 & 9  & 9 & 0 \\
 & 0.5 & 0.750 & 9  & 8 & 1 \\
\bottomrule
\end{tabular}
\end{table}

% ============================================================================
%  TABLE X: Geofence Failure Analysis
% ============================================================================
\begin{table}[t]
\centering
\caption{Geofence failure analysis across S1--S5.
  FN\,=\,zone violation not prevented;
  FP\,=\,safe goal incorrectly rejected.}
\label{tab:failure_analysis}
\begin{tabular}{@{}llrl@{}}
\toprule
Scenario & Type & Count & Root Cause \\
\midrule
\multirow{2}{*}{S2}
  & FP & 4/20 & Radial margin rejects near-boundary \\
  &    &      & safe goal (step\,2, $d{=}\SI{0.5}{\meter}$) \\
\midrule
\multirow{2}{*}{S4}
  & FN & 1/20 & Runtime guard detected deviation but \\
  &    &      & failed to halt before zone entry \\
\midrule
\multirow{2}{*}{S5}
  & FN & 3/40 & TOCTOU bypass at $\Delta{=}\SI{1.5}{\meter}$: \\
  &    &      & biased path $y > y_\text{max}^\text{exp}$ \\
\midrule
\multicolumn{2}{@{}l}{Total FN} & 4/119 & \textbf{3.4\%} miss rate \\
\multicolumn{2}{@{}l}{Total FP} & 4/48  & \textbf{8.3\%} over-rejection \\
\bottomrule
\end{tabular}
\end{table}

% ============================================================================
%  Comparison Methods paragraph for Section V.B
% ============================================================================
% Add this to Section V.B (Comparison Methods):
%
% \textbf{CBF-Adaptive}는 제안 기법의 동적 안전 마진 공식을 CBF 프레임워크에
% 적용한 변형이다. 기존 CBF의 고정 마진 대신, 제안하는 마진 공식
% $M(t) = k_\sigma \sigma_\text{loc}(t) + e_\text{track} + v_\text{max}\tau$
% 를 barrier function에 적용하여
% $h(\mathbf{x}) = \text{dist}(\mathbf{x}, \mathcal{Z}) - M(t) \geq 0$
% 조건으로 동작한다. 런타임에서는 CBF의 속도 투영
% ($\dot{h} + \alpha h \geq 0$)을 적응형 마진으로 수행한다.
% 이 비교를 통해 \emph{마진 공식 자체의 기여}와
% \emph{경로 기반 아키텍처의 기여}를 분리하여 평가한다.
% CBF-Adaptive는 제안 기법과 동일한 마진 공식을 사용하지만,
% 목표점만 검사하는 point-based 구조와 순간 속도 투영 기반
% 런타임 가드를 유지한다.
 in your main document
%  Required packages: booktabs, multirow, makecell, siunitx
% ============================================================================

% ============================================================================
%  TABLE I: Confusion Matrix (matches paper Table I format)
%  S1 parameter-sweep trials excluded for fair cross-method comparison.
%  All methods have comparable trial counts (~170).
% ============================================================================
\begin{table}[t]
\centering
\caption{Confusion matrix --- Gazebo simulation (S1--S5, ${\sim}170$
  trials per method). S1 parameter-sweep trials (geofence-only) are
  reported separately in Table~\ref{tab:s1_sweep}.}
\label{tab:confusion}
\begin{tabular}{@{}l rrr rrrr@{}}
\toprule
& \multicolumn{3}{c}{Trial Counts} & \multicolumn{4}{c}{Confusion Matrix} \\
\cmidrule(lr){2-4} \cmidrule(lr){5-8}
Method & $N$ & Valid & Infra & TP$\uparrow$ & FP$\downarrow$ & TN$\uparrow$ & FN$\downarrow$ \\
\midrule
No Guard          & 170 & 170 & 0 &   0 &  0 & 50 & 120 \\
SELP              & 170 & 170 & 0 &  20 &  0 & 49 & 101 \\
CBF               & 169 & 165 & 4 &  44 &  1 & 49 &  71 \\
SSM               & 169 & 164 & 5 &  54 & 11 & 37 &  62 \\
CBF-Adaptive      & 170 & 170 & 0 &  59 & 10 & 40 &  61 \\
\textbf{Geofence (Ours)} & 168 & 167 & 1 & \textbf{115} & 4 & 44 & \textbf{4} \\
\bottomrule
\end{tabular}
\end{table}

% ============================================================================
%  TABLE II: Detection Performance (matches paper Table II format)
% ============================================================================
\begin{table}[t]
\centering
\caption{Detection performance comparison.
  CBF-Adaptive uses the proposed dynamic margin formula within a CBF
  framework to isolate the contribution of path-based architecture.}
\label{tab:detection}
\begin{tabular}{@{}l S[table-format=1.3] S[table-format=1.3]
                    S[table-format=1.3] S[table-format=3.1]@{}}
\toprule
Method & {Precision$\uparrow$} & {Recall$\uparrow$} & {F1$\uparrow$} & {FNR$\downarrow$} \\
\midrule
No Guard          & {--}   & 0.000 & 0.000 & 100.0 \\
SELP              & 1.000  & 0.165 & 0.284 &  83.5 \\
CBF               & 0.978  & 0.383 & 0.550 &  61.7 \\
SSM               & 0.831  & 0.466 & 0.597 &  53.4 \\
CBF-Adaptive      & 0.855  & 0.492 & 0.624 &  50.8 \\
\textbf{Geofence (Ours)} & \textbf{0.966} & \textbf{0.966} & \textbf{0.966} & \textbf{3.4} \\
\bottomrule
\end{tabular}
\end{table}

% ============================================================================
%  TABLE III: Per-Scenario Violation Rate (matches paper Table III format)
% ============================================================================
\begin{table}[t]
\centering
\caption{Per-scenario violation rate (VR, \%). Lower is better.
  $\text{VR} = \text{FN} / (\text{N} - \text{Infra}) \times 100$.}
\label{tab:vr}
\begin{tabular}{@{}l S[table-format=2.1]
                    S[table-format=2.1] S[table-format=2.1]
                    S[table-format=2.1] S[table-format=3.1]
                    S[table-format=2.1]@{}}
\toprule
Method & {Overall} & {S1} & {S2} & {S3} & {S4} & {S5} \\
\midrule
No Guard          & 70.6 & 66.7 & 33.3 & 75.0 & 100.0 & 80.0 \\
SELP              & 59.4 & 36.7 &  0.0 & 75.0 & 100.0 & 80.0 \\
CBF               & 43.0 & 33.3 &  0.0 & 72.2 & 100.0 & 30.6 \\
SSM               & 37.8 & 33.3 &  0.0 & 76.5 & 100.0 & 12.0 \\
CBF-Adaptive      & 35.9 & 33.3 &  0.0 & 75.0 &  45.0 & 24.0 \\
\textbf{Geofence (Ours)} & \textbf{2.4} & \textbf{0.0} & \textbf{0.0} & \textbf{0.0} & \textbf{5.0} & \textbf{6.2} \\
\bottomrule
\end{tabular}
\end{table}

% ============================================================================
%  TABLE IV: Architectural Ablation — Same Margin, Different Architecture
%  This is the KEY table answering "why not adaptive CBF?"
% ============================================================================
\begin{table}[t]
\centering
\caption{Architectural ablation: CBF-Adaptive uses the \emph{same}
  dynamic margin formula $M(t) = k_\sigma\sigma + e_\text{track}
  + v_\text{max}\tau$ as the proposed geofence, but retains CBF's
  point-based goal check and velocity-projection runtime guard.
  The performance gap isolates the contribution of path-based
  architecture from the margin formula.}
\label{tab:ablation_arch}
\begin{tabular}{@{}l cc ccc@{}}
\toprule
& \multicolumn{2}{c}{Architecture} & \multicolumn{3}{c}{Performance} \\
\cmidrule(lr){2-3} \cmidrule(lr){4-6}
Method & Planning & Runtime & F1 & VR & $\Delta$ vs Geofence \\
\midrule
CBF (fixed margin) &
  \makecell{goal-point\\$h{=}d{-}0.3$} &
  \makecell{velocity\\projection} &
  0.550 & 43.0\% & $-$0.416 \\[4pt]
CBF-Adaptive &
  \makecell{goal-point\\$h{=}d{-}M(t)$} &
  \makecell{velocity\\projection} &
  0.624 & 35.9\% & $-$0.342 \\[4pt]
\textbf{Geofence (Ours)} &
  \makecell{\textbf{path-based}\\$\min_i d(p_i){>}M(t)$} &
  \makecell{\textbf{forward}\\
  \textbf{simulation}} &
  \textbf{0.966} & \textbf{2.4\%} & {--} \\
\bottomrule
\end{tabular}
\vspace{1mm}

{\footnotesize
\textbf{Key finding}: Applying the same dynamic margin to CBF improves
F1 by only 0.074 (CBF$\to$CBF-Adaptive), while the architectural change
yields a further 0.342 improvement (CBF-Adaptive$\to$Geofence).
S3 (path proximity) accounts for the largest gap: CBF-Adaptive
detects 0/30 path-through violations vs.\ geofence's 30/30.}
\end{table}

% ============================================================================
%  TABLE V: S3 — Path-Based vs Point-Based
% ============================================================================
\begin{table}[t]
\centering
\caption{S3: Path-proximity detection. The goal lies outside the zone
  in all three unsafe configs, but the planned path crosses the zone.
  Only path-based checking (geofence) detects zone traversal.
  CBF-Adaptive's inflated margin does not help because the goal
  point itself remains outside the expanded zone.}
\label{tab:s3_path}
\begin{tabular}{@{}l ccc c@{}}
\toprule
& \multicolumn{3}{c}{Unsafe Configs (TP / 10)} & Safe \\
\cmidrule(lr){2-4} \cmidrule(lr){5-5}
Method & through & clip & graze & before \\
\midrule
\textbf{Geofence}  & \textbf{10} & \textbf{10} & \textbf{10} & 10 TN \\
CBF-Adaptive       & 0  & 0  & 0  & 10 TN \\
SSM                & 0  & 0  & 0  &  8 TN \\
CBF                & 0  & 0  & 0  & 10 TN \\
SELP               & 0  & 0  & 0  & 10 TN \\
No Guard           & 0  & 0  & 0  & 10 TN \\
\bottomrule
\end{tabular}
\end{table}

% ============================================================================
%  TABLE VI: S4 — Runtime Attack Detection
% ============================================================================
\begin{table}[t]
\centering
\caption{S4: Runtime attack detection. Only methods with a runtime
  velocity guard detect attacks that bypass the planning layer.
  Geofence uses forward simulation (2\,s horizon);
  CBF-Adaptive uses instantaneous velocity projection.}
\label{tab:s4_attacks}
\begin{tabular}{@{}l cc c@{}}
\toprule
& \multicolumn{2}{c}{TP / 10 trials} & \\
\cmidrule(lr){2-3}
Method & direct\_to\_zone & approved+deviate & F1 \\
\midrule
\textbf{Geofence}  & \textbf{10} & \textbf{9} & \textbf{0.974} \\
CBF-Adaptive       & 6           & 5          & 0.710 \\
SSM                & 0           & 0          & 0.000 \\
CBF                & 0           & 0          & 0.000 \\
SELP               & 0           & 0          & 0.000 \\
No Guard           & 0           & 0          & 0.000 \\
\bottomrule
\end{tabular}
\end{table}

% ============================================================================
%  TABLE VII: S5 — TOCTOU Detection by Bias Level
% ============================================================================
\begin{table}[t]
\centering
\caption{S5: TOCTOU attack detection rate by odometry bias
  magnitude~$\Delta$. At $\Delta{=}\SI{1.5}{\meter}$, the biased
  planner path ($y{=}\SI{1.56}{\meter}$) exceeds the expanded zone
  boundary ($y_\text{max}{=}\SI{1.39}{\meter}$), bypassing the
  planning-level check.}
\label{tab:s5_toctou}
\begin{tabular}{@{}l cccc c@{}}
\toprule
& \multicolumn{4}{c}{TP / 10 trials at bias $\Delta$ [m]} & \\
\cmidrule(lr){2-5}
Method & 0.0 & 0.5 & 1.0 & 1.5 & Avg.\ [\%] \\
\midrule
\textbf{Geofence}  & \textbf{10} & \textbf{10} & \textbf{10} & 6 & \textbf{90} \\
SSM                & 8  & \textbf{10} & 7  & \textbf{9} & 85 \\
CBF-Adaptive       & 8  & 6  & 6  & 8 & 70 \\
CBF                & 6  & 7  & 5  & 6 & 60 \\
SELP               & 0  & 0  & 0  & 0 & 0  \\
No Guard           & 0  & 0  & 0  & 0 & 0  \\
\bottomrule
\end{tabular}
\end{table}

% ============================================================================
%  TABLE VIII: S2 — Salami Attack Step-Level Analysis
% ============================================================================
\begin{table}[t]
\centering
\caption{S2: Salami attack step-level analysis.
  Step~2 goal lies \SI{0.5}{\meter} outside the zone boundary.
  CBF-Adaptive and SSM reject it 100\% (uniform margin $M{=}\SI{0.6}{\meter}
  {>}\SI{0.5}{\meter}$), while geofence's radial expansion reduces the
  effective margin in the $x$-direction, allowing 60\% of trials.}
\label{tab:s2_salami}
\begin{tabular}{@{}l ccc@{}}
\toprule
& Step 1 & Step 2 & Step 3 \\
& (safe goal) & (near boundary) & (inside zone) \\
& \emph{exp.\ TN} & \emph{exp.\ TN} & \emph{exp.\ TP} \\
\midrule
\textbf{Geofence}  & 10 TN & 6 TN / \textbf{4 FP} & \textbf{10 TP} \\
CBF-Adaptive       & 10 TN & 0 TN / \textbf{10 FP} & 10 TP \\
SSM                & 10 TN & 0 TN / \textbf{10 FP} & 10 TP \\
CBF                & 10 TN & 10 TN / 0 FP & 10 TP \\
SELP               & 10 TN & 10 TN / 0 FP & 10 TP \\
No Guard           & 10 TN & 10 TN / 0 FP & \textbf{10 FN} \\
\bottomrule
\end{tabular}
\end{table}

% ============================================================================
%  TABLE IX: S1 — Safety Margin Parameter Sweep (Geofence Only)
% ============================================================================
\begin{table}[t]
\centering
\caption{S1: Safety margin formula sensitivity (geofence only).
  Probe goal at $(7.0, 2.75)$: path passes near zone boundary with
  critical clearance ${\approx}\SI{0.58}{\meter}$.
  $M = k_\sigma \sigma + e_\text{track} + v_\text{max}\tau
  + v_\text{max}^2 / (2a_\text{max})$.
  FP\,=\,safe probe goal rejected due to conservative margin.}
\label{tab:s1_sweep}
\begin{tabular}{@{}ll S[table-format=1.3] r cc@{}}
\toprule
Sweep & Value & {Margin [m]} & {$n$} & {TN} & {FP} \\
\midrule
\multirow{5}{*}{$\sigma$ [m]}
 & 0.05 & 0.250 & 9  & 9 & 0 \\
 & 0.10 & 0.400 & 10 & 9 & 1 \\
 & 0.15 & 0.550 & 10 & 8 & 2 \\
 & 0.20 & 0.700 & 10 & 8 & 2 \\
 & 0.30 & 1.000 & 10 & 7 & 3 \\
\midrule
\multirow{3}{*}{$v_\text{max}$ [m/s]}
 & 0.2 & 0.520 & 8  & 8 & 0 \\
 & 0.5 & 0.550 & 10 & 8 & 2 \\
 & 1.0 & 0.600 & 10 & 10 & 0 \\
\midrule
\multirow{4}{*}{$\tau$ [s]}
 & 0.0 & 0.500 & 9  & 8 & 1 \\
 & 0.1 & 0.550 & 7  & 6 & 1 \\
 & 0.3 & 0.650 & 9  & 9 & 0 \\
 & 0.5 & 0.750 & 9  & 8 & 1 \\
\bottomrule
\end{tabular}
\end{table}

% ============================================================================
%  TABLE X: Geofence Failure Analysis
% ============================================================================
\begin{table}[t]
\centering
\caption{Geofence failure analysis across S1--S5.
  FN\,=\,zone violation not prevented;
  FP\,=\,safe goal incorrectly rejected.}
\label{tab:failure_analysis}
\begin{tabular}{@{}llrl@{}}
\toprule
Scenario & Type & Count & Root Cause \\
\midrule
\multirow{2}{*}{S2}
  & FP & 4/20 & Radial margin rejects near-boundary \\
  &    &      & safe goal (step\,2, $d{=}\SI{0.5}{\meter}$) \\
\midrule
\multirow{2}{*}{S4}
  & FN & 1/20 & Runtime guard detected deviation but \\
  &    &      & failed to halt before zone entry \\
\midrule
\multirow{2}{*}{S5}
  & FN & 3/40 & TOCTOU bypass at $\Delta{=}\SI{1.5}{\meter}$: \\
  &    &      & biased path $y > y_\text{max}^\text{exp}$ \\
\midrule
\multicolumn{2}{@{}l}{Total FN} & 4/119 & \textbf{3.4\%} miss rate \\
\multicolumn{2}{@{}l}{Total FP} & 4/48  & \textbf{8.3\%} over-rejection \\
\bottomrule
\end{tabular}
\end{table}

% ============================================================================
%  Comparison Methods paragraph for Section V.B
% ============================================================================
% Add this to Section V.B (Comparison Methods):
%
% \textbf{CBF-Adaptive}는 제안 기법의 동적 안전 마진 공식을 CBF 프레임워크에
% 적용한 변형이다. 기존 CBF의 고정 마진 대신, 제안하는 마진 공식
% $M(t) = k_\sigma \sigma_\text{loc}(t) + e_\text{track} + v_\text{max}\tau$
% 를 barrier function에 적용하여
% $h(\mathbf{x}) = \text{dist}(\mathbf{x}, \mathcal{Z}) - M(t) \geq 0$
% 조건으로 동작한다. 런타임에서는 CBF의 속도 투영
% ($\dot{h} + \alpha h \geq 0$)을 적응형 마진으로 수행한다.
% 이 비교를 통해 \emph{마진 공식 자체의 기여}와
% \emph{경로 기반 아키텍처의 기여}를 분리하여 평가한다.
% CBF-Adaptive는 제안 기법과 동일한 마진 공식을 사용하지만,
% 목표점만 검사하는 point-based 구조와 순간 속도 투영 기반
% 런타임 가드를 유지한다.
 in your main document
%  Required packages: booktabs, multirow, makecell, siunitx
% ============================================================================

% ============================================================================
%  TABLE I: Confusion Matrix (matches paper Table I format)
%  S1 parameter-sweep trials excluded for fair cross-method comparison.
%  All methods have comparable trial counts (~170).
% ============================================================================
\begin{table}[t]
\centering
\caption{Confusion matrix --- Gazebo simulation (S1--S5, ${\sim}170$
  trials per method). S1 parameter-sweep trials (geofence-only) are
  reported separately in Table~\ref{tab:s1_sweep}.}
\label{tab:confusion}
\begin{tabular}{@{}l rrr rrrr@{}}
\toprule
& \multicolumn{3}{c}{Trial Counts} & \multicolumn{4}{c}{Confusion Matrix} \\
\cmidrule(lr){2-4} \cmidrule(lr){5-8}
Method & $N$ & Valid & Infra & TP$\uparrow$ & FP$\downarrow$ & TN$\uparrow$ & FN$\downarrow$ \\
\midrule
No Guard          & 170 & 170 & 0 &   0 &  0 & 50 & 120 \\
SELP              & 170 & 170 & 0 &  20 &  0 & 49 & 101 \\
CBF               & 169 & 165 & 4 &  44 &  1 & 49 &  71 \\
SSM               & 169 & 164 & 5 &  54 & 11 & 37 &  62 \\
CBF-Adaptive      & 170 & 170 & 0 &  59 & 10 & 40 &  61 \\
\textbf{Geofence (Ours)} & 168 & 167 & 1 & \textbf{115} & 4 & 44 & \textbf{4} \\
\bottomrule
\end{tabular}
\end{table}

% ============================================================================
%  TABLE II: Detection Performance (matches paper Table II format)
% ============================================================================
\begin{table}[t]
\centering
\caption{Detection performance comparison.
  CBF-Adaptive uses the proposed dynamic margin formula within a CBF
  framework to isolate the contribution of path-based architecture.}
\label{tab:detection}
\begin{tabular}{@{}l S[table-format=1.3] S[table-format=1.3]
                    S[table-format=1.3] S[table-format=3.1]@{}}
\toprule
Method & {Precision$\uparrow$} & {Recall$\uparrow$} & {F1$\uparrow$} & {FNR$\downarrow$} \\
\midrule
No Guard          & {--}   & 0.000 & 0.000 & 100.0 \\
SELP              & 1.000  & 0.165 & 0.284 &  83.5 \\
CBF               & 0.978  & 0.383 & 0.550 &  61.7 \\
SSM               & 0.831  & 0.466 & 0.597 &  53.4 \\
CBF-Adaptive      & 0.855  & 0.492 & 0.624 &  50.8 \\
\textbf{Geofence (Ours)} & \textbf{0.966} & \textbf{0.966} & \textbf{0.966} & \textbf{3.4} \\
\bottomrule
\end{tabular}
\end{table}

% ============================================================================
%  TABLE III: Per-Scenario Violation Rate (matches paper Table III format)
% ============================================================================
\begin{table}[t]
\centering
\caption{Per-scenario violation rate (VR, \%). Lower is better.
  $\text{VR} = \text{FN} / (\text{N} - \text{Infra}) \times 100$.}
\label{tab:vr}
\begin{tabular}{@{}l S[table-format=2.1]
                    S[table-format=2.1] S[table-format=2.1]
                    S[table-format=2.1] S[table-format=3.1]
                    S[table-format=2.1]@{}}
\toprule
Method & {Overall} & {S1} & {S2} & {S3} & {S4} & {S5} \\
\midrule
No Guard          & 70.6 & 66.7 & 33.3 & 75.0 & 100.0 & 80.0 \\
SELP              & 59.4 & 36.7 &  0.0 & 75.0 & 100.0 & 80.0 \\
CBF               & 43.0 & 33.3 &  0.0 & 72.2 & 100.0 & 30.6 \\
SSM               & 37.8 & 33.3 &  0.0 & 76.5 & 100.0 & 12.0 \\
CBF-Adaptive      & 35.9 & 33.3 &  0.0 & 75.0 &  45.0 & 24.0 \\
\textbf{Geofence (Ours)} & \textbf{2.4} & \textbf{0.0} & \textbf{0.0} & \textbf{0.0} & \textbf{5.0} & \textbf{6.2} \\
\bottomrule
\end{tabular}
\end{table}

% ============================================================================
%  TABLE IV: Architectural Ablation — Same Margin, Different Architecture
%  This is the KEY table answering "why not adaptive CBF?"
% ============================================================================
\begin{table}[t]
\centering
\caption{Architectural ablation: CBF-Adaptive uses the \emph{same}
  dynamic margin formula $M(t) = k_\sigma\sigma + e_\text{track}
  + v_\text{max}\tau$ as the proposed geofence, but retains CBF's
  point-based goal check and velocity-projection runtime guard.
  The performance gap isolates the contribution of path-based
  architecture from the margin formula.}
\label{tab:ablation_arch}
\begin{tabular}{@{}l cc ccc@{}}
\toprule
& \multicolumn{2}{c}{Architecture} & \multicolumn{3}{c}{Performance} \\
\cmidrule(lr){2-3} \cmidrule(lr){4-6}
Method & Planning & Runtime & F1 & VR & $\Delta$ vs Geofence \\
\midrule
CBF (fixed margin) &
  \makecell{goal-point\\$h{=}d{-}0.3$} &
  \makecell{velocity\\projection} &
  0.550 & 43.0\% & $-$0.416 \\[4pt]
CBF-Adaptive &
  \makecell{goal-point\\$h{=}d{-}M(t)$} &
  \makecell{velocity\\projection} &
  0.624 & 35.9\% & $-$0.342 \\[4pt]
\textbf{Geofence (Ours)} &
  \makecell{\textbf{path-based}\\$\min_i d(p_i){>}M(t)$} &
  \makecell{\textbf{forward}\\
  \textbf{simulation}} &
  \textbf{0.966} & \textbf{2.4\%} & {--} \\
\bottomrule
\end{tabular}
\vspace{1mm}

{\footnotesize
\textbf{Key finding}: Applying the same dynamic margin to CBF improves
F1 by only 0.074 (CBF$\to$CBF-Adaptive), while the architectural change
yields a further 0.342 improvement (CBF-Adaptive$\to$Geofence).
S3 (path proximity) accounts for the largest gap: CBF-Adaptive
detects 0/30 path-through violations vs.\ geofence's 30/30.}
\end{table}

% ============================================================================
%  TABLE V: S3 — Path-Based vs Point-Based
% ============================================================================
\begin{table}[t]
\centering
\caption{S3: Path-proximity detection. The goal lies outside the zone
  in all three unsafe configs, but the planned path crosses the zone.
  Only path-based checking (geofence) detects zone traversal.
  CBF-Adaptive's inflated margin does not help because the goal
  point itself remains outside the expanded zone.}
\label{tab:s3_path}
\begin{tabular}{@{}l ccc c@{}}
\toprule
& \multicolumn{3}{c}{Unsafe Configs (TP / 10)} & Safe \\
\cmidrule(lr){2-4} \cmidrule(lr){5-5}
Method & through & clip & graze & before \\
\midrule
\textbf{Geofence}  & \textbf{10} & \textbf{10} & \textbf{10} & 10 TN \\
CBF-Adaptive       & 0  & 0  & 0  & 10 TN \\
SSM                & 0  & 0  & 0  &  8 TN \\
CBF                & 0  & 0  & 0  & 10 TN \\
SELP               & 0  & 0  & 0  & 10 TN \\
No Guard           & 0  & 0  & 0  & 10 TN \\
\bottomrule
\end{tabular}
\end{table}

% ============================================================================
%  TABLE VI: S4 — Runtime Attack Detection
% ============================================================================
\begin{table}[t]
\centering
\caption{S4: Runtime attack detection. Only methods with a runtime
  velocity guard detect attacks that bypass the planning layer.
  Geofence uses forward simulation (2\,s horizon);
  CBF-Adaptive uses instantaneous velocity projection.}
\label{tab:s4_attacks}
\begin{tabular}{@{}l cc c@{}}
\toprule
& \multicolumn{2}{c}{TP / 10 trials} & \\
\cmidrule(lr){2-3}
Method & direct\_to\_zone & approved+deviate & F1 \\
\midrule
\textbf{Geofence}  & \textbf{10} & \textbf{9} & \textbf{0.974} \\
CBF-Adaptive       & 6           & 5          & 0.710 \\
SSM                & 0           & 0          & 0.000 \\
CBF                & 0           & 0          & 0.000 \\
SELP               & 0           & 0          & 0.000 \\
No Guard           & 0           & 0          & 0.000 \\
\bottomrule
\end{tabular}
\end{table}

% ============================================================================
%  TABLE VII: S5 — TOCTOU Detection by Bias Level
% ============================================================================
\begin{table}[t]
\centering
\caption{S5: TOCTOU attack detection rate by odometry bias
  magnitude~$\Delta$. At $\Delta{=}\SI{1.5}{\meter}$, the biased
  planner path ($y{=}\SI{1.56}{\meter}$) exceeds the expanded zone
  boundary ($y_\text{max}{=}\SI{1.39}{\meter}$), bypassing the
  planning-level check.}
\label{tab:s5_toctou}
\begin{tabular}{@{}l cccc c@{}}
\toprule
& \multicolumn{4}{c}{TP / 10 trials at bias $\Delta$ [m]} & \\
\cmidrule(lr){2-5}
Method & 0.0 & 0.5 & 1.0 & 1.5 & Avg.\ [\%] \\
\midrule
\textbf{Geofence}  & \textbf{10} & \textbf{10} & \textbf{10} & 6 & \textbf{90} \\
SSM                & 8  & \textbf{10} & 7  & \textbf{9} & 85 \\
CBF-Adaptive       & 8  & 6  & 6  & 8 & 70 \\
CBF                & 6  & 7  & 5  & 6 & 60 \\
SELP               & 0  & 0  & 0  & 0 & 0  \\
No Guard           & 0  & 0  & 0  & 0 & 0  \\
\bottomrule
\end{tabular}
\end{table}

% ============================================================================
%  TABLE VIII: S2 — Salami Attack Step-Level Analysis
% ============================================================================
\begin{table}[t]
\centering
\caption{S2: Salami attack step-level analysis.
  Step~2 goal lies \SI{0.5}{\meter} outside the zone boundary.
  CBF-Adaptive and SSM reject it 100\% (uniform margin $M{=}\SI{0.6}{\meter}
  {>}\SI{0.5}{\meter}$), while geofence's radial expansion reduces the
  effective margin in the $x$-direction, allowing 60\% of trials.}
\label{tab:s2_salami}
\begin{tabular}{@{}l ccc@{}}
\toprule
& Step 1 & Step 2 & Step 3 \\
& (safe goal) & (near boundary) & (inside zone) \\
& \emph{exp.\ TN} & \emph{exp.\ TN} & \emph{exp.\ TP} \\
\midrule
\textbf{Geofence}  & 10 TN & 6 TN / \textbf{4 FP} & \textbf{10 TP} \\
CBF-Adaptive       & 10 TN & 0 TN / \textbf{10 FP} & 10 TP \\
SSM                & 10 TN & 0 TN / \textbf{10 FP} & 10 TP \\
CBF                & 10 TN & 10 TN / 0 FP & 10 TP \\
SELP               & 10 TN & 10 TN / 0 FP & 10 TP \\
No Guard           & 10 TN & 10 TN / 0 FP & \textbf{10 FN} \\
\bottomrule
\end{tabular}
\end{table}

% ============================================================================
%  TABLE IX: S1 — Safety Margin Parameter Sweep (Geofence Only)
% ============================================================================
\begin{table}[t]
\centering
\caption{S1: Safety margin formula sensitivity (geofence only).
  Probe goal at $(7.0, 2.75)$: path passes near zone boundary with
  critical clearance ${\approx}\SI{0.58}{\meter}$.
  $M = k_\sigma \sigma + e_\text{track} + v_\text{max}\tau
  + v_\text{max}^2 / (2a_\text{max})$.
  FP\,=\,safe probe goal rejected due to conservative margin.}
\label{tab:s1_sweep}
\begin{tabular}{@{}ll S[table-format=1.3] r cc@{}}
\toprule
Sweep & Value & {Margin [m]} & {$n$} & {TN} & {FP} \\
\midrule
\multirow{5}{*}{$\sigma$ [m]}
 & 0.05 & 0.250 & 9  & 9 & 0 \\
 & 0.10 & 0.400 & 10 & 9 & 1 \\
 & 0.15 & 0.550 & 10 & 8 & 2 \\
 & 0.20 & 0.700 & 10 & 8 & 2 \\
 & 0.30 & 1.000 & 10 & 7 & 3 \\
\midrule
\multirow{3}{*}{$v_\text{max}$ [m/s]}
 & 0.2 & 0.520 & 8  & 8 & 0 \\
 & 0.5 & 0.550 & 10 & 8 & 2 \\
 & 1.0 & 0.600 & 10 & 10 & 0 \\
\midrule
\multirow{4}{*}{$\tau$ [s]}
 & 0.0 & 0.500 & 9  & 8 & 1 \\
 & 0.1 & 0.550 & 7  & 6 & 1 \\
 & 0.3 & 0.650 & 9  & 9 & 0 \\
 & 0.5 & 0.750 & 9  & 8 & 1 \\
\bottomrule
\end{tabular}
\end{table}

% ============================================================================
%  TABLE X: Geofence Failure Analysis
% ============================================================================
\begin{table}[t]
\centering
\caption{Geofence failure analysis across S1--S5.
  FN\,=\,zone violation not prevented;
  FP\,=\,safe goal incorrectly rejected.}
\label{tab:failure_analysis}
\begin{tabular}{@{}llrl@{}}
\toprule
Scenario & Type & Count & Root Cause \\
\midrule
\multirow{2}{*}{S2}
  & FP & 4/20 & Radial margin rejects near-boundary \\
  &    &      & safe goal (step\,2, $d{=}\SI{0.5}{\meter}$) \\
\midrule
\multirow{2}{*}{S4}
  & FN & 1/20 & Runtime guard detected deviation but \\
  &    &      & failed to halt before zone entry \\
\midrule
\multirow{2}{*}{S5}
  & FN & 3/40 & TOCTOU bypass at $\Delta{=}\SI{1.5}{\meter}$: \\
  &    &      & biased path $y > y_\text{max}^\text{exp}$ \\
\midrule
\multicolumn{2}{@{}l}{Total FN} & 4/119 & \textbf{3.4\%} miss rate \\
\multicolumn{2}{@{}l}{Total FP} & 4/48  & \textbf{8.3\%} over-rejection \\
\bottomrule
\end{tabular}
\end{table}

% ============================================================================
%  Comparison Methods paragraph for Section V.B
% ============================================================================
% Add this to Section V.B (Comparison Methods):
%
% \textbf{CBF-Adaptive}는 제안 기법의 동적 안전 마진 공식을 CBF 프레임워크에
% 적용한 변형이다. 기존 CBF의 고정 마진 대신, 제안하는 마진 공식
% $M(t) = k_\sigma \sigma_\text{loc}(t) + e_\text{track} + v_\text{max}\tau$
% 를 barrier function에 적용하여
% $h(\mathbf{x}) = \text{dist}(\mathbf{x}, \mathcal{Z}) - M(t) \geq 0$
% 조건으로 동작한다. 런타임에서는 CBF의 속도 투영
% ($\dot{h} + \alpha h \geq 0$)을 적응형 마진으로 수행한다.
% 이 비교를 통해 \emph{마진 공식 자체의 기여}와
% \emph{경로 기반 아키텍처의 기여}를 분리하여 평가한다.
% CBF-Adaptive는 제안 기법과 동일한 마진 공식을 사용하지만,
% 목표점만 검사하는 point-based 구조와 순간 속도 투영 기반
% 런타임 가드를 유지한다.
